\textbf{COLUNI 2025 -- Questão 11}\par\medskip Nos Jogos Olímpicos de Tóquio 2020, o Time Brasil superou o recorde de medalhas conquistadas (21) em uma única edição; igualou o número de ouros obtidos no Rio 2016 (7), subiu ao pódio em 13 modalidades, um feito inédito, e terminou em 12° lugar no quadro de medalhas (7 ouros, 6 pratas e 8 bronzes).\par

\noindent Disponível em: https://www.cob.org.br/pt/cob/time-brasil/brasil-nosjogos/medalhas-olimpicas. Acesso em: 08 jul. 2024.

\par

O Brasil participou de 23 Jogos Olímpicos ao longo de sua história e conquistou medalhas em 20 dessas edições. O quadro abaixo apresenta as medalhas conquistadas.

\begin{center}
\begin{tabular}{|l|c|c|c|c|}
\hline
\textbf{EDIÇÃO} & \shortstack{\textbf{MEDALHAS}\\\textbf{OURO}} & \shortstack{\textbf{MEDALHAS}\\\textbf{PRATA}} & \shortstack{\textbf{MEDALHAS}\\\textbf{BRONZE}} & \shortstack{\textbf{MEDALHAS}\\\textbf{TOTAL}} \\
\hline
Tóquio 2020 & 7 & 6 & 8 & 21 \\
Rio 2016 & 7 & 6 & 6 & 19 \\
Londres 2012 & 3 & 5 & 9 & 17 \\
Pequim 2008 & 3 & 4 & 10 & 17 \\
Atenas 2006 & 5 & 2 & 3 & 10 \\
Sydney 2000 & 0 & 6 & 6 & 12 \\
Atlanta 1996 & 3 & 3 & 9 & 15 \\
Barcelona 1992 & 2 & 1 & 0 & 3 \\
Seul 1988 & 1 & 2 & 3 & 6 \\
Los Angeles 1984 & 1 & 5 & 2 & 8 \\
Moscou 1980 & 2 & 0 & 2 & 4 \\
Montreal 1976 & 0 & 0 & 2 & 2 \\
Munique 1972 & 0 & 0 & 2 & 2 \\
Cidade do México 1968 & 0 & 1 & 2 & 3 \\
Tóquio 1964 & 0 & 0 & 1 & 1 \\
Roma 1960 & 0 & 0 & 2 & 2 \\
Melbourne 1956 & 1 & 0 & 0 & 1 \\
Helsinque 1952 & 1 & 0 & 2 & 3 \\
Londres 1948 & 0 & 0 & 1 & 1 \\
Antuérpia 1920 & 1 & 1 & 1 & 3 \\
\textbf{TOTAL} & \textbf{37} & \textbf{42} & \textbf{71} & \textbf{150} \\
\hline
\end{tabular}
\end{center}

\noindent Disponível em: https://www.cob.org.br/pt/cob/time-brasil/brasil-nosjogos/medalhas-olimpicas. Acesso em: 08 jul. 2024.

\par

Considerando a tabela acima, calcule:

\begin{enumerate}
\item[I.] Na edição de Atlanta, a percentagem de medalhas de bronze conquistadas em relação ao total.
\item[II.] A média de medalhas de prata conquistadas ao longo dessas 20 edições.
\end{enumerate}

Marque a alternativa que apresenta a resposta CORRETA para os itens I e II, respectivamente.
\begin{enumerate}
\item 60\% e 2,1.
\item 60\% e 3,1.
\item 63\% e 3,5.
\item 63\% e 2,5.
\end{enumerate}
